%Glossareinträge ab hier
\newglossaryentry{workaround}
{
    name=Workaround,
    description={Ein Workaround ist, ein Plan oder eine Methode zur Umgehung eines Problems (z.B. im Bereich der Computersoftware), ohne es zu beseitigen}
}
\newglossaryentry{known error}
{
    name=Known Error,
    description={Ein Known Error ist ein Problem, für das es eine dokumentierte Ursache und einen Workaround gibt}
}
\newglossaryentry{fifo}
{
    name=FIFO,
    description={Die Kurzbezeichnung für First-in-first-out; Ein Prioritätsprinzip (Priorität) der Warteschlangentheorie, nach dem zuerst ankommende Transaktionen zuerst bedient werden, welche \ac{ua} bei der Reihenfolgeplanung angewandt wird}
}
\newglossaryentry{agile}
{
    name=Agile,
    description={Ein Oberbegriff für eine Sammlung an Frameworks und Techniken, die es Teams und Einzelpersonen ermöglichen, auf eine Weise zu arbeiten, die durch Zusammenarbeit, Priorisierung, iterative und inkrementelle Bereitstellung und Timeboxing gekennzeichnet ist. Es gibt mehrere spezifische Methoden wie Scrum, Lean und Kanban, die als agil eingestuft werden.}
}
\newglossaryentry{lean}
{
    name=Lean,
    description={Ein Ansatz, der sich auf die Verbesserung von Arbeitsabläufen konzentriert, indem der Wert durch die Vermeidung von Verschwendung maximiert wird.}
}
\newglossaryentry{devops}
{
    name=DevOps,
    description={Eine Organisationskultur, die darauf abzeilt, den Wertfluss zum Kunden zu verbsessern. DevOps konzentiriert sich aud CALMS: Culture, Automation, Lean, Measurement und Sharing}
}
\newglossaryentry{effizienz}
{
    name=Effizienz,
    description={Misst, wie gut Ressourcen genutzt werden, um ein Ziel zu erreichen. Es geht um das „Wie“ - weniger Ressourcenverbrauch gilt als effizienter.}
}
\newglossaryentry{effektivitaet}
{
    name=Effektivität,
    description={Bezieht sich darauf, inwieweit die gewünschten Ergebnisse oder Ziele erreicht werden. Es geht um das „Was“ - das Erreichen des beabsichtigten Ergebnisses.}
}
\newglossaryentry{carveout}
{
	name=Carve-Out,
	description={Abspaltung bzw. Veräußerung von Unternehmensteilen (z.B. einzelne Geschäftsbereiche), die i.d.R. als Bestandteil eines Restrukturierungsplans innerhalb des veräußernden Unternehmens einzuordnen ist. Unternehmen veräußern einzelne Unternehmensteile häufig im Rahmen der Fokussierung auf das Kerngeschäft, das mit Hilfe der Verkaufserlöse noch ausgebaut werden kann. Insofern umfasst der Begriff Carve-out sowohl die Veräußerung von rentablen als auch unrentablen Unternehmensteilen.\footcite{Dr.SteffenHundtDCO}}
}
\newglossaryentry{cmdb}
{
	name=CMDB,
	description={Die Abkürzung "CMDB" steht für "Configuration Management Database". Sie bezeichnet eine Datenbank mit allen relevanten Informationen zur IT-Infrastruktur. In der CMDB werden Objekte (CIs - Configuration Items) mit zusätzlichen Informationen erfasst, die Beziehungen zwischen diesen Objekten dargestellt (Standort, Integration, einer Person zugeordnet / übergeben), Objekte mit Dokumenten verknüpft, die verschiedenen Zustände der Objekte (betriebsbereit / defekt / in Reparatur) hinterlegt und die verantwortlichen Personen dokumentiert. Da diese Datenbank bereits alle vorhandenen Informationen enthält und durch die  Prozesse des Konfigurationsmanagements kontinuierlich gepflegt wird, dient die CMDB als Basis für weitere wichtige Managementsysteme wie ITSM oder ISMS. Die CMDB setzt die Komponenten zueinander in Beziehung und bietet so einen Überblick über die Zusammenhänge und Abhängigkeiten aller Komponenten über den gesamten Lebenszyklus hinweg}
}
\newglossaryentry{heyjo}
{
	name=Hey-Jo-Prinzip,
	description={Das Hey-Joe-Prinzip beschreibt die Auswirkungen der im IT-Support problematischen „Nachbarschaftshilfe” (Peer Support), die dabei den vorgesehenen Arbeitsprozess (Workflow) umgeht. Das geschieht, indem der Anwender die Hilfe für eine Problemlösung über eine inoffizielle Anfrage in seinem firmeninternen Bekanntenkreis erfragt, anstelle den eigentlich dafür vorgesehenen Service Desk zu kontaktieren, was plakativ mit “Hey Joe” beschrieben wird}
}
\newglossaryentry{adhoc}
{
	name={ad hoc},
	description={Die lateinische Phrase „ad hoc“ meint im Deutschen meist „aus dem Augenblick heraus“ und beschreibt damit etwa spontane Handlungen, Aussagen, Urteile usw. Die Redewendung „aus dem Stegreif heraus“ ist ebenfalls ein Synonym für ad hoc}
}
\newglossaryentry{bab}
{
	name={Betriebsabrechnungsbogen},
	description={Der Betriebsabrechnungsbogen (BAB) dient der internen Kostenverrechnung als Kalkulationsschema. Mit seiner Hilfe werden die Einzelkosten erfasst und die Gemeinkosten auf die Kostenstellen verteilt. Für das Unternehmen wird so ersichtlich, wo welche Kosten entstanden sind. Es kann Zuschlagssätze für die Selbstkostenkalkulation ermitteln und erhält darüber hinaus eine Grundlage, um die Gemeinkosten zu analysieren und zu steuern.}
}